\metatitle{Constructions and examples}
\metadescription{Allowed-color words and their matching, orientation, path,
and matrix models; lifts of ordinary matroids and weighted assignments.}

\section[constructions]{Constructions and examples}

We give two constructions: words with allowed colors and prescribed quotas,
and lifts of arbitrary ordinary matroids. Matchings, orientations, interval
paths, and the matrix presentation below describe the first construction.
A weighted example at the end illustrates optimization by cyclic exchanges.

\section[allowedColors]{The allowed-color construction}

Fix quotas $r_1,\dotsc,r_k$. At each position $e$, choose allowed colors
$L_e\subseteq[k]$, and take every word with $u_e\in L_e$ and the prescribed
quotas.

\begin{theorem}[Allowed-color construction]
Every nonempty family obtained in this way is a colored matroid.
\end{theorem}

\begin{proof}
For two words $u,v$, draw an arrow $u_e\to v_e$ at each disagreement.
The resulting directed multigraph is balanced because the quotas agree.
Every arrow lies on a directed simple cycle. Switch one through the
prescribed position. The quotas are preserved, and each new entry was
already allowed at that position in the other word.
\end{proof}

\subsection[matchingAssignments]{Matchings}

Join position $e$ to color $c$ when $c\in L_e$. A word chooses one
incident edge at each position and exactly $r_c$ edges at color $c$.
Equivalently, replace color $c$ by $r_c$ labeled slots with identical
neighbors. A perfect matching assigns positions to slots; forgetting slot
labels gives the word.

The ordinary \defin{transversal matroid} records which position sets can
be matched to distinct slots \cite{Oxley2011}. Here we fill every quota
and retain the colors of the slots.

\subsection[graphOrientations]{Graph orientations}

When each list consists of the endpoints of an edge, the positions form
an edge set and the colors a vertex set. Orient each edge toward its
assigned color. Quotas prescribe indegrees, so the words are all
orientations with those indegrees. Cyclic exchange reverses a directed
cycle. The \href{graphic.htm}{graph chapter} gives the corresponding cut
criterion and circuits.

\subsection[latticePaths]{Interval paths}

Order $E=[n]$, choose an interval $I_c$ for each color, and require
$B(u,c)\subseteq I_c$. For each color, the permitted sets are the bases
of a uniform rank-$r_c$ matroid on $I_c$, with loops elsewhere.
This is an ordinary lattice-path matroid: the north-step positions can
be any $r_c$ positions in the interval. Taking all compatible partitions
is exactly the allowed-color construction. For ordinary lattice-path
matroids, see \cite{BoninMierNoy2003}.

Prefix-count bounds give a larger construction; see
\href{lattice-paths.htm}{colored lattice-path matroids}.

\subsection[matrixPresentations]{Matrix presentations}

\begin{proposition}[Matrix realization of allowed colors]
Every allowed-color colored matroid has matrices $X_c$ over $\mathbb R$ such
that its words are exactly the compatible partitions satisfying, for every $c$,
\[
 \det X_c[:,B(u,c)]\ne0.
\]
\end{proposition}

\begin{proof}
For color $c$, put zero columns at positions where $c$ is forbidden. On the
allowed positions, use distinct columns of an $r_c$-row Vandermonde matrix.
Its bases are exactly the $r_c$-subsets of the allowed positions. Hence the
compatible basis partitions are precisely the allowed-color words, which
satisfy symmetric cycle exchange by the allowed-color theorem.
\end{proof}

\begin{example}[One family with interval and matrix presentations]
Take quotas $(2,1,1)$ on four positions. Allow color $1$ everywhere,
color $2$ on $\{1,2,3\}$, and color $3$ on $\{1,2,4\}$. This is the
\href{independence.htm#mixedCircuits}{seven-word family}. Its matrix
presentation is
\[
 \begin{gathered}
 X_1=\begin{pmatrix}1&1&1&1\\0&1&2&3\end{pmatrix},\\
 X_2=\begin{pmatrix}1&1&1&0\end{pmatrix},\\
 X_3=\begin{pmatrix}1&1&0&1\end{pmatrix}.
 \end{gathered}
\]
Every two columns of $X_1$ are independent; $X_2,X_3$ allow exactly
the stated singleton choices. Ordering the positions as $3,1,2,4$
also gives an interval presentation: color $2$ is allowed in the
first three positions and color $3$ in the last three.
\end{example}

\section[representedLifts]{Lifts of ordinary matroids}

We now start with any ordinary matroid. Let $N$
be a rank-$r$ matroid on $[n]$, choose a prime $p\geq n$, and choose distinct
residues $s_1,\dotsc,s_n$ modulo $p$. Take colors
\[
 \ell_0,\dotsc,\ell_{p-1},\rho_0,\dotsc,\rho_{p-1}.
\]
For each $j\in[n]$, introduce the $p$ positions
\[
 M_j=\{(\ell_i,\rho_{i+s_j}):i\in\mathbb Z/p\mathbb Z\}.
\]
These are pairwise disjoint perfect matchings between the $\ell$-colors and
the $\rho$-colors.
For a basis $B$ of $N$, color every position in $M_j$ by its $\ell$-endpoint
when $j\in B$, and by its $\rho$-endpoint when $j\notin B$.

\begin{theorem}[Matroid lift]
The words obtained from the bases of $N$ form a colored matroid. Every
$\ell$-color has quota $r$, and every $\rho$-color has quota $n-r$.
If $N$ is representable, then every color projection of the lifted colored
matroid is representable.
\end{theorem}

\begin{proof}
Each matching $M_j$ contains every $\ell$-color and every $\rho$-color once,
which gives the stated quotas. Take two bases $B,C$ and a prescribed
disagreement in $M_j$, where after interchanging the parents we may assume
$j\in B\setminus C$. Symmetric basis exchange gives
$h\in C\setminus B$ such that
\[
 B\setminus\{j\}\cup\{h\}
 \qquad\text{and}\qquad
 C\setminus\{h\}\cup\{j\}
\]
are bases of $N$.

The union of $M_j$ and $M_h$ is a cycle through all $2p$ colors. Indeed,
following the two matchings changes the index of an $\ell$-color by the
nonzero residue $s_j-s_h$, which generates $\mathbb Z/p\mathbb Z$. The
disagreement arrows on this cycle alternate between the two matchings.
Switching the cycle gives exactly the words associated with the two exchanged
bases above, and the cycle contains the prescribed position.

For a fixed $\ell_i$, its selected positions are indexed by the basis $B$ of
$N$. For a fixed $\rho_i$, they are indexed by the complementary basis of
$N^*$. Thus the color projections are relabeled copies of $N$ and $N^*$,
with loops at nonincident positions. Representability is preserved by duality.
\end{proof}

\begin{corollary}[Long exchanges and weighted marginals]
In the matroid lift, every cyclic exchange with both resulting words
in the family has length $2p$. Every weighted completion polynomial is
a positive scalar times a positive diagonal substitution of the basis polynomial
of $N$ or $N^*$. Hence all of them are Lorentzian.
\end{corollary}

\begin{proof}
Each matching $M_j$ receives either all its $\ell$-endpoints or all its
$\rho$-endpoints. An exchange with resulting words in the family must
switch whole matchings, at least one in each direction. Thus at least
$2p$ positions change. A cyclic exchange uses each color at most once,
giving equality.

Distinguish color $\ell_i$, and write $x_j$ for its variable on $M_j$.
Let $a_j$ be the product of the weights at the other $\ell$-endpoints
of $M_j$, and $b_j$ the product at all its $\rho$-endpoints. With
\[
 b_N(\mathbf x)=\sum_{B\in\mathcal B(N)}\prod_{j\in B}x_j,
\]
the completion polynomial is
\[
 \left(\prod_{j=1}^n b_j\right)
 b_N((a_jx_j/b_j)_{j=1}^n).
\]
For color $\rho_i$, interchange the two choices and use the
complementary basis, giving the same formula with $N^*$.
Matroid basis polynomials are Lorentzian, and positive diagonal
substitutions preserve Lorentzianity
\cite[Theorems~2.10 and 3.10]{BrandenHuh2020}.
\end{proof}

For example, $N=U_{3,6}$ and $p=7$ give twenty words on forty-two
positions with fourteen colors, each of quota three. Every cyclic
exchange with resulting words in the family uses all fourteen
colors. Thus the weighted Lorentzian property permits arbitrarily
large minimum exchange length.

\begin{example}[A graphic and representable lift]
Take $N=U_{2,3}$, the rank-two graphic matroid of a triangle, and use the
three shift matchings of $K_{3,3}$. The construction gives three words on
nine positions and six colors. Each $\ell$-color has quota $2$, each
$\rho$-color has quota $1$, and an exchange between two words switches the
six-color cycle formed by the corresponding pair of matchings.
\end{example}

\section[weightedAssignments]{Weighted assignments}

The \href{exchange-graphs.htm#scoreOptimization}{optimality criterion}
says that improving cyclic exchanges lead to an optimum. Greedy acceptance
of high-scoring assignments can stop short of it, as the following example
shows.

\begin{example}[Eight positions with quotas $(3,3,2)$]
We assign colors $1,2,3$ to eight positions, using the colors three, three,
and two times, respectively. The table gives the score of each allowed
assignment; a dash means that the color is forbidden.
\[
\begin{array}{c|rrr}
 \text{position}&\text{Color }1&\text{Color }2&\text{Color }3\\ \hline
 1&20&14&-\\
 2&19&-&13\\
 3&-&11&10\\
 4&18&12&-\\
 5&17&9&-\\
 6&16&8&-\\
 7&15&-&7\\
 8&-&5&6
\end{array}
\]
We take all words respecting the quotas and allowed colors. This is a
colored matroid by the allowed-color proposition. Write $a(w)$ for the sum
of the eight assignment scores.

Greedy scans the allowed position--color pairs in decreasing score order,
accepting a pair whenever the partial assignment still extends to a feasible
word. It first accepts color $1$ at positions $1,2,4$, with scores $20,19,18$.
Color $1$ is now full, so it rejects the scores $17,16,15$. It completes to
\[
 u=\mathtt{11212233},\qquad a(u)=98,
\]
with color classes $\{1,2,4\}$, $\{3,5,6\}$, and $\{7,8\}$.

Compare $u$ with $v=\mathtt{23321112}$. Positions $1,3,7$ give the cycle
$1\to2\to3\to1$. Switching these positions in both words gives
\[
\begin{array}{c|cccccccc|r}
 &1&2&3&4&5&6&7&8&\text{score}\\ \hline
 u&\mathbf1&1&\mathbf2&1&2&2&\mathbf3&3&98\\
 v&\mathbf2&3&\mathbf3&2&1&1&\mathbf1&2&102\\ \hline
 u'&\mathbf2&1&\mathbf3&1&2&2&\mathbf1&3&99\\
 v'&\mathbf1&3&\mathbf2&2&1&1&\mathbf3&2&101
\end{array}
\]
Both resulting words are feasible, and their total score is unchanged.
Continuing from $u'$ gives the improving path
\[
\begin{array}{c|c|r}
 \text{word}&\text{positions changed}&\text{score}\\ \hline
 \mathtt{11212233}&-&98\\
 \mathtt{21312213}&1,3,7&99\\
 \mathtt{23211213}&2,3,5&102\\
 \mathtt{23221113}&4,6&104
\end{array}
\]
The first two moves are three-color cycles; the last is a two-color swap.

The final word is optimal. Subtract $7$ from every allowed score in color
$1$. Every word loses exactly $21$, and the adjusted row maxima occur at
colors $2,3,2,2,1,1,1,3$. These form the feasible word
$\mathtt{23221113}$, so it attains the sum of all adjusted row maxima and
maximizes the original score.
\end{example}
